

\documentclass[aspectratio=169,8pt]{beamer}
\usetheme{Madrid}
\usepackage[utf8]{inputenc}
\usepackage[T1]{fontenc}
\usepackage{amsmath,amssymb}
\usepackage{booktabs}
\usepackage{enumitem}
\setlist[enumerate]{label=\Alph*.,leftmargin=*,itemsep=1pt,topsep=2pt}
\setlist[itemize]{leftmargin=*,itemsep=1pt,topsep=2pt}
\setbeamerfont{block body}{size=\tiny}
\setbeamerfont{block title}{size=\scriptsize}
\setbeamerfont{frametitle}{size=\footnotesize}
\setbeamerfont{normal text}{size=\tiny}
\setbeamertemplate{navigation symbols}{}
\setbeamertemplate{itemize item}{\tiny$\bullet$}
\setbeamertemplate{itemize subitem}{\tiny$\circ$}

\title[GARP RAI]{GARP Risk \& AI --- Q351--Q450}
\subtitle{Unofficial Exam Prep}
\author{Unofficial}
\date{\today}

\begin{document}

\begin{frame}{Q560 --- Distributive Justice}
	\begin{block}{Question}
		Distributive justice is primarily concerned with:
		\begin{enumerate}
			\item Fair allocation of benefits and burdens
			\item Fair legal procedures
			\item Protection of privacy
			\item Model explainability
		\end{enumerate}
	\end{block}
	
	\begin{exampleblock}{Answer: A}\end{exampleblock}
	
	\begin{block}{Explanation}
		Distributive justice concerns how resources, opportunities, benefits, and burdens are allocated.
	\end{block}
\end{frame}

\begin{frame}{Q561 --- Procedural Justice}
	\begin{block}{Question}
		Procedural justice is MOST concerned with:
		\begin{enumerate}
			\item Equal outcomes
			\item Fairness of decision-making processes
			\item Utility maximization
			\item Resource redistribution
		\end{enumerate}
	\end{block}
	
	\begin{exampleblock}{Answer: B}\end{exampleblock}
\end{frame}

\begin{frame}{Q562 --- Procedural Justice Requirement}
	\begin{block}{Question}
		Which of the following is generally required for procedural justice?
		\begin{enumerate}
			\item Impartial application of rules
			\item Equal outcomes
			\item Merit-based rewards
			\item Utility maximization
		\end{enumerate}
	\end{block}
	
	\begin{exampleblock}{Answer: A}\end{exampleblock}
\end{frame}

\begin{frame}{Q563 --- Justice Conflict}
	\begin{block}{Question}
		Which statement regarding justice is correct?
		\begin{enumerate}
			\item Different theories of distributive justice always produce identical outcomes
			\item Distributive justice concepts may conflict with one another
			\item Procedural justice eliminates distributive justice concerns
			\item Social justice eliminates autonomy concerns
		\end{enumerate}
	\end{block}
	
	\begin{exampleblock}{Answer: B}\end{exampleblock}
\end{frame}

\begin{frame}{Q564 --- Egalitarianism}
	\begin{block}{Question}
		Egalitarianism is MOST closely associated with:
		\begin{enumerate}
			\item Equal treatment or equal distribution
			\item Rewarding achievement
			\item Maximizing aggregate welfare
			\item Maximizing shareholder value
		\end{enumerate}
	\end{block}
	
	\begin{exampleblock}{Answer: A}\end{exampleblock}
\end{frame}

\begin{frame}{Q565 --- Utilitarianism}
	\begin{block}{Question}
		A utilitarian approach generally seeks to:
		\begin{enumerate}
			\item Benefit the greatest number of people
			\item Provide identical outcomes
			\item Reward merit only
			\item Eliminate all inequality
		\end{enumerate}
	\end{block}
	
	\begin{exampleblock}{Answer: A}\end{exampleblock}
\end{frame}

\begin{frame}{Q566 --- Meritocracy}
	\begin{block}{Question}
		Under a meritocratic approach, benefits are allocated primarily according to:
		\begin{enumerate}
			\item Need
			\item Luck
			\item Achievement or contribution
			\item Equal outcomes
		\end{enumerate}
	\end{block}
	
	\begin{exampleblock}{Answer: C}\end{exampleblock}
\end{frame}

\begin{frame}{Q567 --- Meritocracy Definition Trap}
	\begin{block}{Question}
		Which statement is TRUE regarding meritocracy?
		\begin{enumerate}
			\item Meritocracy is not considered a distributive justice principle
			\item Meritocracy allocates benefits according to achievement
			\item Meritocracy requires equal outcomes
			\item Meritocracy requires equal wealth
		\end{enumerate}
	\end{block}
	
	\begin{exampleblock}{Answer: B}\end{exampleblock}
\end{frame}

\begin{frame}{Q568 --- Autonomy and Justice}
	\begin{block}{Question}
		According to the Responsible AI module, social justice generally:
		\begin{enumerate}
			\item Requires human autonomy
			\item Eliminates the need for autonomy
			\item Is incompatible with autonomy
			\item Requires machine autonomy instead of human autonomy
		\end{enumerate}
	\end{block}
	
	\begin{exampleblock}{Answer: A}\end{exampleblock}
\end{frame}

\begin{frame}{Q569 --- Justice in AI}
	\begin{block}{Question}
		An AI hiring model systematically disadvantages a qualified demographic group. Which ethical principle is MOST directly implicated?
		\begin{enumerate}
			\item Nonmaleficence
			\item Justice
			\item Explainability
			\item Security
		\end{enumerate}
	\end{block}
	
	\begin{exampleblock}{Answer: B}\end{exampleblock}
\end{frame}

\begin{frame}{Q570 --- Conflicting Justice Concepts}
	\begin{block}{Question}
		An algorithm distributes resources equally among all citizens. A utilitarian analysis suggests resources should be concentrated among a smaller group to maximize overall welfare. This illustrates:
		\begin{enumerate}
			\item Procedural failure
			\item Explainability failure
			\item Conflict between justice concepts
			\item Violation of autonomy
		\end{enumerate}
	\end{block}
	
	\begin{exampleblock}{Answer: C}\end{exampleblock}
\end{frame}

\end{document}